\begin{python0}
from solutions import *; clear()
\end{python0}
\sectionthree{Function call graph}

The diagram above that describes function call
\begin{python}
from latextool_basic import *
p = Plot()

p += Rect(0, 8, 3, 9, label=r'\texttt{f(0)}', linewidth=0.05)
p += Rect(0, 6, 3, 7, label=r'\texttt{f(1)}', linewidth=0.05)
p += Rect(0, 4, 3, 5, label=r'\texttt{f(2)}', linewidth=0.05)
p += Rect(0, 2, 3, 3, label=r'\texttt{f(3)}', linewidth=0.05)
p += Rect(0, 0, 3, 1, label=r'\texttt{main}', linewidth=0.05)
#up
zaphod = 1
for i in range(4):
    zaphod += 1
    p += Line(points=[(0.75, zaphod - 1), (0.75, zaphod)], linewidth=0.075, endstyle='>')
    zaphod += 1

#down
for i in range(4):
    zaphod -= 1
    p += Line(points=[(2.25, zaphod), (2.25, zaphod - 1)], linewidth=0.075, endstyle='>')
    zaphod -= 1
print(p)
\end{python}

is called a \EMPHASIZE{function call graph}. Usually a function call graph
is used to describe the calling, so the return path is usually not
shown. If I'm only interested in the recursive calls say of f(3), then I
would not include the box for main(). So the function call graph for
f(3) looks like this:

\begin{python}
from latextool_basic import *
p = Plot()

p += Rect(0, 8, 3, 9, label=r'\texttt{f(0)}', linewidth=0.05)
p += Rect(0, 6, 3, 7, label=r'\texttt{f(1)}', linewidth=0.05)
p += Rect(0, 4, 3, 5, label=r'\texttt{f(2)}', linewidth=0.05)
p += Rect(0, 2, 3, 3, label=r'\texttt{f(3)}', linewidth=0.05)
#up
zaphod = 3
for i in range(3):
    zaphod += 1
    p += Line(points=[(1.5, zaphod - 1), (1.5, zaphod)], linewidth=0.075, endstyle='>')
    zaphod += 1

print(p)
\end{python}

is important because, if there are no loops in the function, then the
number of boxes in the function call graph more or less denotes the
amount of work that has to be done to compute f(3). For instance,
omitting the box for \texttt{main()}, f(3) involves 4 boxes.


